\section{Project Scope and Assumptions}
\label{sec:project_scope}
\textbf{Scope:} The project is strictly scoped to building a research prototype. Complex cryptographic proof generation is primarily modeled via a mock delay, and execution is limited to simplified transfer state transitions. 

\textbf{Assumptions:} We assume that the local dev-chain (Hardhat) environment accurately reflects the relative (though not absolute) L1 execution costs of mainnet. We also assume that the synthetic Poisson workload provides a sufficient baseline to measure relative queueing dynamics, acknowledging that real-world DeFi traffic is significantly more bursty.

\subsection{Research Question Mapping}
Table \ref{tab:rq_mapping} maps the research questions to the objectives, components, and report sections.

\begin{table}[H]
\centering
\caption{Research Question Mapping}
\label{tab:rq_mapping}
\compacttable
\begin{tabularx}{\textwidth}{@{}p{0.10\textwidth}Y Y Y p{0.16\textwidth}@{}}
\toprule
\textbf{RQ} & \textbf{Related objective} & \textbf{Related component} & \textbf{Related metric} & \textbf{Report section} \\
\midrule
RQ1 & Evaluate batch size effects & Sequencer, Executor, Submitter & TPS, latency, gas cost, proof delay & Sections 3.9, 5.7, 5.8 \\
RQ2 & Evaluate sequencing policies & Sequencer, Executor & Fairness, latency, execution correctness & Sections 3.9, 5.10 \\
RQ3 & Compare data availability modes & Submitter, L1 contracts & DA footprint and cost per transaction & Sections 3.9, 5.11 \\
RQ4 & Identify practical bottlenecks & Integration pipeline & Development issues and trade-offs & Sections 4.13, 5.17 \\
\bottomrule
\end{tabularx}
\end{table}
\normalsize
